\section*{Bab \ref{ch:g12:reaction-rates} --- Laju Reaksi}

\begin{solution}{exo:g12:reaction-rates:1}
Suhu; luas bidang sentuh; konsentrasi; suhu.
\end{solution}

\begin{solution}{exo:g12:reaction-rates:2}
\qty{13.9}{min} (biru) dan \qty{6.9}{min} (merah). Setelah dua kali \omterm{def:g12:reaction-rates:half-life}{waktu
paruh}, $10.0 / 4 = \qty{2.5}{mmol/L}$.
\end{solution}

\begin{solution}{exo:g12:reaction-rates:3}
Kemiringan $(5.0 - 1.0)/10 = \qty{0.40}{mmol.L^{-1}.min^{-1}}$: itulah
\omterm{def:g12:reaction-rates:rate}{laju pembentukan} pada \qty{5}{min}.
\end{solution}

\begin{solution}{exo:g12:reaction-rates:4}
Pendinginan dan \omterm{def:g10:concentration-and-dilution:dilution}{pengenceran} membuat reaksinya begitu lambat sehingga
komposisinya tidak lagi berubah selama \omterm{def:g11:titration:titration}{titrasi}. Komposisi yang terukur
adalah komposisi \omterm{def:g6:pure-substances-and-mixtures:mixture}{campuran} saat reaksinya dihentikan, yaitu saat
penuangan.
\end{solution}

\begin{solution}{exo:g12:reaction-rates:5}
\qty{30}{min} adalah 3 kali \omterm{def:g12:reaction-rates:half-life}{waktu paruh}: tersisa $1/8$. \qty{1}{h}
adalah 6 kali \omterm{def:g12:reaction-rates:half-life}{waktu paruh}: $1/64$, sekitar \qty{1.6}{\percent}.
\end{solution}

\begin{solution}{exo:g12:reaction-rates:6}
$\ln[\ce{A}]$: 2.079, 1.733, 1.386, 1.040, 0.693. Nilainya turun 0.346
setiap \qty{5}{min}: sebuah garis lurus, orde satu. $k = 0.346/5 =
\qty{0.0693}{min^{-1}}$ dan $t_{1/2} = \ln 2 / k = \qty{10.0}{min}$
(seperti yang terlihat langsung: 8.00, 4.00, 2.00 setiap
\qty{10}{min}).
\end{solution}

\begin{solution}{exo:g12:reaction-rates:7}
$k = 0.693 / 25 = \qty{0.0277}{s^{-1}}$.
\end{solution}

\begin{solution}{exo:g12:reaction-rates:8}
$0.100\,e^{-0.20} = \qty{0.0819}{mol/L}$; $0.100\,e^{-0.60} =
\qty{0.0549}{mol/L}$; $0.100\,e^{-2.0} = \qty{0.0135}{mol/L}$.
\end{solution}

\begin{solution}{exo:g12:reaction-rates:9}
Pada \qty{13.9}{min}, $[\ce{A}] = \qty{5.0}{mmol/L}$; kemiringan garis
singgungnya mendekati $-0.25$, dan $-k[\ce{A}] = -0.050 \times 5.0 =
\qty{-0.25}{mmol.L^{-1}.min^{-1}}$: setengah nilai awalnya.
\end{solution}

\begin{solution}{exo:g12:reaction-rates:10}
Iodin adalah satu-satunya spesies berwarna: menurut hukum Beer--Lambert
$A = \varepsilon \ell [\ce{I2}]$. \omterm{def:g12:reaction-rates:rate}{Laju pembentukan} adalah kemiringan
garis singgung kurva $[\ce{I2}](t)$, yaitu kemiringan garis singgung
$A(t)$ dibagi dengan $\varepsilon \ell$.
\end{solution}

\begin{solution}{exo:g12:reaction-rates:11}
Laju awal $k[\ce{A}]_0$ dua kali lebih besar pada percobaan kedua; \omterm{def:g12:reaction-rates:half-life}{waktu
paruhnya} sama, $\ln 2 / k$, tidak bergantung pada $[\ce{A}]_0$.
\end{solution}

\begin{solution}{exo:g12:reaction-rates:12}
$t = \ln(1/f)/k$. Dalam kali \omterm{def:g12:reaction-rates:half-life}{waktu paruh}: $\ln(1/f)/\ln 2$. Untuk
\qty{1}{\percent}, $\ln 100 / \ln 2 = 6.6$; untuk \qty{0.1}{\percent},
$\ln 1000 / \ln 2 =
10.0$.
\end{solution}

\begin{solution}{exo:g12:reaction-rates:13}
\qty{90}{\percent} terpakai, $f = 0.10$: $t = \ln 10 / k$, yaitu
\qty{46}{min} pada \qty{20}{\celsius} dan \qty{23}{min} pada
\qty{30}{\celsius}. \omterm{def:g12:reaction-rates:first-order}{Tetapan lajunya} menjadi dua kali lipat dalam
\qty{10}{\celsius}: patokan itu berlaku di sini.
\end{solution}

\begin{solution}{exo:g12:reaction-rates:14}
Suhunya turun sampai mendekati \qty{0}{\celsius}, dan konsentrasinya
terbagi 10. Dengan laju yang sebanding dengan hasil kali dua
konsentrasi, \omterm{def:g10:concentration-and-dilution:dilution}{pengenceran} saja membaginya dengan $10 \times 10 = 100$.
\end{solution}

\begin{solution}{exo:g12:reaction-rates:15}
$v(t) = -\dfrac{d}{dt}\big([\ce{A}]_0 e^{-kt}\big) = k[\ce{A}]_0 e^{-kt}$.
Laju awal $k[\ce{A}]_0$; lajunya menjadi setengah ketika
$e^{-kt} = 1/2$, yaitu setelah satu kali \omterm{def:g12:reaction-rates:half-life}{waktu paruh}.
\end{solution}

\begin{solution}{pb:g12:reaction-rates:1}
\textbf{1.} Menurut hukum Beer--Lambert, karena zat warna itu
satu-satunya spesies yang menyerap pada panjang gelombang ini;
larutannya harus cukup encer ($A$ di bawah sekitar 1).

\textbf{2.} Produk itu tidak menyerap cahaya tampak: produk itu tidak
menambah apa pun pada $A$.

\textbf{3.} Agar konsentrasinya hampir tidak berubah: lajunya lalu
hanya bergantung pada zat warna.

\textbf{4.} $A = 0$ (blanko).

\textbf{5.} $A$ turun dari 0.800 menjadi 0.402 dalam \qty{6}{min}:
sekitar \qty{6}{min}.

\textbf{6.} Dari 0.402 pada \qty{6}{min} menjadi 0.199 pada
\qty{12}{min}: terbagi dua lagi dalam \qty{6}{min}.

\textbf{7.} $-0.223$, $-0.448$, $-0.691$, $-0.911$, $-1.155$, $-1.366$,
$-1.614$, $-1.945$, $-2.551$, $-3.079$.

\textbf{8.} Titik-titiknya terletak pada sebuah garis lurus: reaksinya
berorde satu terhadap zat warna.

\textbf{9.} $(-2.551 + 0.223)/20 = \qty{-0.116}{min^{-1}}$.

\textbf{10.} $k = \qty{0.116}{min^{-1}}$.

\textbf{11.} $0.693 / 0.116 = \qty{6.0}{min}$, seperti yang terbaca pada
tabel.

\textbf{12.} $k A_0 = 0.116 \times 0.800 = \qty{0.093}{min^{-1}}$ dalam
satuan \omterm{def:g11:absorbance:absorbance}{absorbans} per menit.

\textbf{13.} Sama, \qty{6.0}{min}: \omterm{def:g12:reaction-rates:half-life}{waktu paruh} orde satu tidak
bergantung pada nilai awalnya.

\textbf{14.} $0.800\, e^{-0.116 \times 30} = 0.025$.

\textbf{15.} $t = \ln 100 / k = 4.61 / 0.116 = \qty{40}{min}$.

\textbf{16.} $40 / 6.0 = 6.6$ kali \omterm{def:g12:reaction-rates:half-life}{waktu paruh}.

\textbf{17.} $t = \ln(0.800/0.010)/k = 4.38 / 0.116 = \qty{38}{min}$.

\textbf{18.} Tetap 6.6 kali \omterm{def:g12:reaction-rates:half-life}{waktu paruh}, kini $6.6 \times 3.0 = \qty{20}{min}$.

\textbf{19.} Sekitar 6.6 kali \omterm{def:g12:reaction-rates:half-life}{waktu paruh}, yaitu sekitar \qty{40}{min}
di sini.
\end{solution}
